\documentclass[12pt,a4paper]{article}

%=========================================================
%                   ENCODAGE ET POLICES
%=========================================================
\usepackage[utf8]{inputenc}
\usepackage[T1]{fontenc}
\usepackage[french]{babel}
\usepackage{lmodern}

%=========================================================
%                     MATHÉMATIQUES
%=========================================================
\usepackage{amsmath,amssymb,amsfonts,mathrsfs}
\usepackage{tkz-tab}
\everymath{\displaystyle}

%=========================================================
%                     MISE EN PAGE
%=========================================================
\usepackage[left=1.5cm,right=1.5cm,top=1.5cm,bottom=1.8cm]{geometry}
\usepackage{multicol}
\usepackage{float}
\usepackage{array,multirow}
\usepackage{enumitem}
\usepackage{color, soul}

%=========================================================
%                GRAPHISMES ET COULEURS
%=========================================================
\usepackage{tikz}
\usetikzlibrary{arrows.meta,calc}
\usepackage{pgfplots}
\pgfplotsset{compat=1.15}

\usepackage[most]{tcolorbox}
\tcbuselibrary{skins,hooks}

%=========================================================
%                  LIENS ET ARRIÈRE-PLAN
%=========================================================
\usepackage[colorlinks=true,linkcolor=blue,urlcolor=blue,citecolor=blue]{hyperref}
\usepackage{eso-pic}

%=========================================================
%                   COULEURS ET STYLES
%=========================================================
\definecolor{myred}{RGB}{180, 40, 40}
\definecolor{myblue}{RGB}{20, 50, 120}
\definecolor{mygreen}{RGB}{30, 120, 40}

%=========================================================
%                       FILIGRANE
%=========================================================
\AddToShipoutPicture{
    \AtTextCenter{
        \makebox[0pt]{
            \rotatebox{45}{
                \textcolor[gray]{0.92}{
                    \fontsize{5cm}{5cm}\selectfont PGB
                }
            }
        }
    }
}

\begin{document}

% En-tête
\begin{center}
    \Large\textbf{\underline{\textcolor{myred}{Racine Carrée et Nombres Réels}}}\\[0.2cm]
    \normalsize\textbf{Prof : M. BA} \hfill \textbf{Classe : 3 ème}\\[0.1cm]
    \textbf{Durée : 6 heures} \hfill \textbf{Année scolaire : 2026 -- 2027}
\end{center}
\hrule
\vspace{0.3cm}

%=========================================================
%    I. Définitions et notations, valeur exacte et valeur approchée
%=========================================================
\section*{\color{myred}I. Définitions et notations, valeur exacte et valeur approchée}


\subsection*{\color{myred}1) Activité}

On considère un carré $ABCD$ de côté $3$ cm.

\begin{enumerate}[label=\textbf{\arabic*.}]
    \item Calculer l'aire du carré $ABCD$.
    \item Déterminer la longueur du côté d'un carré dont l'aire est
    $25\ \text{cm}^2$, puis celle d'un carré dont l'aire est
    $49\ \text{cm}^2$.
    \item Existe-t-il un carré dont l'aire est $10\ \text{cm}^2$ ?
    Si oui, déterminer la longueur exacte de son côté.
    \item À l'aide de la calculatrice, donner une valeur approchée
    de cette longueur au centième près.
    \item Peut-on écrire cette longueur sous la forme d'un nombre
    décimal ayant un nombre fini de chiffres après la virgule ?
\end{enumerate}

\subsection*{\color{myred} 2) Définitions et notations}
\begin{tcolorbox}[colback=red!5!white,colframe=myred,title=]
Soit $a$ un nombre rationnel positif ou nul ($a \ge 0$). La \textbf{racine carrée} de $a$, notée $\mathbf{\sqrt{a}}$, est le nombre positif dont le carré est égal à $a$.\\
L'ensemble des nombres réels est noté $\mathbb{R}$.
\end{tcolorbox}

\noindent \textbf{Conséquence immédiate de la définition :}
Pour tout nombre réel $a \ge 0$, on a : 
\[
(\sqrt{a})^2 = a \quad \text{et} \quad \sqrt{a^2} = a
\]

\noindent \textbf{NB :} 
Le symbole $\sqrt{\phantom{a}}$ est appelé \textit{radical}. La racine carrée d'un nombre négatif n'existe pas dans $\mathbb{R}$ (ex: $\sqrt{-4}$ est impossible).

\subsection*{\color{myred} 3) Valeur exacte et valeur approchée}
\begin{itemize}
    \item \textbf{Valeur exacte :} Soit la racine carrée de $7$, elle se note $\sqrt{7}$. La valeur exacte de ce nombre est égale à $\sqrt{7}$.(ex : $\sqrt{2}$, $\sqrt{3}$, $\frac{2\sqrt{5}}{3}$).
    \item \textbf{Valeur approchée :} $\sqrt{7} = 2{,}64575131\dots$\\ $2{,}645$ est la valeur approchée à $0{,}001$ ($10^{-3}$) près par défaut et $2{,}646$ est la valeur approchée par excès.
\end{itemize}
\subsection*{\color{myred} Exercices d'application}

\noindent \textbf{Exercice 1 :} Soit $B = \sqrt{11}$.
\begin{enumerate}[label=\alph*)]
    \item Calcule $B^2$.
    \item Donne la valeur exacte de $B$.
    \item Détermine la valeur approchée de $B$ à deux unités près par défaut.
    \item Donne la valeur approchée de $B$ à $0{,}01$ près par excès.
    \item Encadre $B$ à $10^{-2}$ près.
\end{enumerate}

\noindent \textbf{Exercice 2 :} Calculer ou simplifier si possible : 
\[
A = \sqrt{121} \quad ; \quad C = (\sqrt{15})^2 \quad ; \quad D = \sqrt{(-7)^2} \quad ; \quad E = \sqrt{-25}
\]

\vspace{0.4cm}
\hrule
\vspace{0.4cm}

\subsection*{\color{myred} Correction des Exercices}

\noindent \textbf{Correction de l'Exercice 1 :}
\begin{enumerate}[label=\alph*)]
    \item $B^2 = (\sqrt{11})^2 = 11$.
    \item La valeur exacte de $B$ est $\mathbf{\sqrt{11}}$.
    \item Sachant que $\sqrt{11} \approx 3{,}3166\dots$, sa valeur approchée à deux unités près par défaut est $\mathbf{3}$.
    \item La valeur approchée de $B$ à $0{,}01$ près par excès est $\mathbf{3{,}32}$.
    \item L'encadrement de $B$ à $10^{-2}$ près ($0{,}01$ près) est : $\mathbf{3{,}31 < \sqrt{11} < 3{,}32}$.
\end{enumerate}

\vspace{0.3cm}
\noindent \textbf{Correction de l'Exercice 2 :}
\begin{itemize}
    \item $\mathbf{A = \sqrt{121} = 11}$ (car $11^2 = 121$).
    \item $\mathbf{C = (\sqrt{15})^2 = 15}$ (par définition, $(\sqrt{a})^2 = a$ pour $a \ge 0$).
    \item $\mathbf{D = \sqrt{(-7)^2} = |-7| = 7}$ (la racine carrée d'un carré est toujours positive, égale à la valeur absolue).
    \item $\mathbf{E = \sqrt{-25}}$ : \textbf{Impossible} (ou n'a pas de sens dans $\mathbb{R}$), car le radicande est négatif ($-25 < 0$).
\end{itemize}

%=========================================================
%    II. Nombres irrationnels : ensemble IR des nombres réels
%=========================================================
\section*{\color{myred}II. Nombres irrationnels : ensemble $\mathbb{R}$ des nombres réels}

\subsection*{\color{myred} 1) Activité}
\begin{enumerate}
    \item A l’aide d’une calculatrice, détermine : $\sqrt{3}$, $\sqrt{29}$, $\sqrt{2}$,$\pi$,$\sqrt{196}$
    \item Parmi ces nombres, quels sont ceux qui ne sont pas des nombres rationnels ?
Ceux qui sont des nombres rationnels ?
\end{enumerate} 

\subsection*{\color{myred} 2) Définitions et notations}
\begin{itemize}
    \item \textbf{Nombres rationnels ($\mathbb{Q}$)} Nombres pouvant s'écrire sous forme de fraction $\frac{p}{q}$ (avec $p \in \mathbb{Z}$ et $q \in \mathbb{N}^*$).
    \item \textbf{Nombres irrationnels :} Nombres qui ne peuvent pas s'écrire sous forme de fraction (leurs écritures décimales sont infinies et non périodiques, ex : $\sqrt{2}$, $\pi$).
    \item \textbf{Ensemble des nombres réels ($\mathbb{R}$) :} L'ensemble qui regroupe tous les nombres rationnels et irrationnels.
\end{itemize}

\noindent \textbf{Relation d'inclusion :}
$
\mathbb{N} \subset \mathbb{Z} \subset \mathbb{D} \subset \mathbb{Q} \subset \mathbb{R}
$

\subsection*{\color{myred} 3) Exercice d’application}
Parmi les nombres suivants, indiquez ceux qui appartiennent à $\mathbb{Q}$ et ceux qui appartiennent aux irrationnels :
\[
-3 \quad ; \quad \frac{5}{2} \quad ; \quad \sqrt{9} \quad ; \quad \sqrt{7} \quad ; \quad \pi
\]

%=========================================================
%    III. Propriétés de la racine carrée
%=========================================================
\section*{\color{myred}III. Propriétés de la racine carrée}

\subsection*{\color{myred} 1) Activité}
Calculer et comparer :
\begin{itemize}
    \item $\sqrt{9} \times \sqrt{16}$ et $\sqrt{9 \times 16}$
   \item $\frac{\sqrt{36}}{\sqrt{4}}$ et $\sqrt{\frac{36}{4}}$
\end{itemize}

\subsection*{\color{myred} 2) Propriétés}
\begin{itemize}
    \item \textbf{Produit :}La racine carrée du produit de deux nombres réels positifs est égale au produit des racines de ces deux réels. 
  $$\sqrt{a \times b} = \sqrt{a} \times \sqrt{b} \quad (\text{pour } a \ge 0, b \ge 0)$$
  \textit{Exemple :} $\sqrt{50} = \sqrt{25 \times 2} = \sqrt{25} \times \sqrt{2} = 5\sqrt{2}$.
    \item \textbf{Quotient :}Le quotient des racines carrées de deux nombres réels positifs est égal à la racine carrée du quotient de ces deux réels. Autrement dit, si $a \in \mathbb{R}_+$ et $b \in \mathbb{R}_*^+$, alors :
  $$\sqrt{\frac{a}{b}} = \frac{\sqrt{a}}{\sqrt{b}}$$
  \textit{Exemples :} $\sqrt{\frac{9}{16}} = \frac{\sqrt{9}}{\sqrt{16}} = \frac{3}{4}$.
\end{itemize}

\begin{tcolorbox}[colback=red!10!white,colframe=myred,title=\textbf{Attention / Erreur fréquente}]
\textbf{La somme et la différence ne se distribuent pas !}
$$\sqrt{a + b} \neq \sqrt{a} + \sqrt{b} \quad \text{et} \quad \sqrt{a - b} \neq \sqrt{a} - \sqrt{b}$$
\textit{Contre-exemple :} $\sqrt{9 + 16} = \sqrt{25} = 5$, alors que $\sqrt{9} + \sqrt{16} = 3 + 4 = 7$. Or $5 \neq 7$.
\end{tcolorbox}

%=========================================================
%    IV. Calculs sur les radicaux
%=========================================================
\section*{\color{myred}IV. Calculs sur les radicaux}

\subsection*{A. Sommes algébriques}

\subsubsection*{\color{myred} 1) Activité}
\begin{enumerate}[label=\alph*)]
    \item Écris les radicaux suivants sous la forme $a\sqrt{b}$ avec $a \in \mathbb{N}$ et $b \in \mathbb{N}$ :
    \[
    \sqrt{18} \quad ; \quad \sqrt{50} \quad ; \quad \sqrt{32} \quad ; \quad \sqrt{48}
    \]
    \item Déduis-en une écriture simplifiée des sommes algébriques suivantes :
    \[
    C = \sqrt{18} + 2\sqrt{50} + 3\sqrt{32} \qquad \text{et} \qquad D = \sqrt{50} - \sqrt{32} + \sqrt{48}
    \]
\end{enumerate}
\subsubsection*{\color{myred} 1) Activité}
Soient les nombres :
\[
F = \frac{5}{\sqrt{2}} \quad ; \quad G = \frac{4}{3\sqrt{7}} \quad ; \quad H = \frac{1}{\sqrt{6} - 2}
\]
Multiplie chacun de ces nombres par le nombre qui convient pour obtenir un dénominateur sans radical.
\subsubsection*{\color{myred} 2) Propriétés}
\begin{itemize}
    \item Quels que soient les réels $b$ et $c$ et le réel positif $a$, on a :
    \[
    b\sqrt{a} + c\sqrt{a} = (b+c)\sqrt{a} \qquad \text{et} \qquad b\sqrt{a} - c\sqrt{a} = (b-c)\sqrt{a}
    \]
\item Si $p \in \mathbb{N}$ et $a \ge 0$, alors :
\begin{itemize}
    \item Pour une puissance paire :
    \[
    \sqrt{a^{2p}} = \sqrt{(a^p)^2} = a^p
    \]
    \item Pour une puissance impaire :
    \[
    \sqrt{a^{2p+1}} = \sqrt{a^{2p} \times a} = \sqrt{a^{2p}} \times \sqrt{a} = a^p\sqrt{a}
    \]
\end{itemize}
\end{itemize}
\subsubsection*{\color{myred} Exemple}
$B = \sqrt{45} - 2\sqrt{5} + \sqrt{20} = 3\sqrt{5} - 2\sqrt{5} + 2\sqrt{5} = 3\sqrt{5}$.

$
E = \sqrt{5} + \sqrt{20} - \sqrt{80} = \sqrt{5} + 2\sqrt{5} - 4\sqrt{5} = (1 + 2 - 4)\sqrt{5} = -\sqrt{5}
$
\subsubsection*{\color{myred} Exemples}
Calculer et simplifier les expressions suivantes :
\[
B = \sqrt{3^{8}} \quad ; \quad C = \sqrt{3^{5}} \quad ; \quad D = \sqrt{3^{11}}
\]

\noindent \textbf{Correction :}
\begin{itemize}
    \item $B = \sqrt{3^{8}} = \sqrt{3^{2 \times 4}} = 3^4 = 81$
    \item $C = \sqrt{3^{5}} = \sqrt{3^{2 \times 2 + 1}} = 3^2\sqrt{3} = 9\sqrt{3}$
    \item $D = \sqrt{3^{11}} = \sqrt{3^{2 \times 5 + 1}} = 3^5\sqrt{3} = 243\sqrt{3}$
\end{itemize}
\subsection*{B. Expressions conjuguées ; rendre rationnel le dénominateur}

\subsubsection*{\color{myred} 1) Activité}
Soient les nombres :
\[
F = \frac{5}{\sqrt{2}} \quad ; \quad G = \frac{4}{3\sqrt{7}} \quad ; \quad H = \frac{1}{\sqrt{6} - 2}
\]
Multiplie chacun de ces nombres par le nombre qui convient pour obtenir un dénominateur sans radical.
\subsubsection*{\color{myred} 2) Vocabulaire}
\begin{itemize}
\item Le conjugué de $(a + \sqrt{b})$ est $(a - \sqrt{b})$.
\item Le conjugué de $(\sqrt{a} - \sqrt{b})$ est $(\sqrt{a} + \sqrt{b})$.
\end{itemize}
\subsubsection*{\color{myred} Exemple}
\begin{itemize}
    \item L'expression conjuguée de $7 + \sqrt{5}$ est $7 - \sqrt{5}$.
    \item L'expression conjuguée de $\sqrt{3} - \sqrt{2}$ est $\sqrt{3} + \sqrt{2}$.
\end{itemize}
\subsubsection*{\color{myred} 3) Méthode et Exemple}
Pour rendre rationnel un dénominateur comportant une somme ou une différence de radicaux, on multiplie le numérateur et le dénominateur par l'expression conjuguée du dénominateur.

\noindent \textit{Exemple :}
\[
\frac{2}{\sqrt{3} - 1} = \frac{2(\sqrt{3} + 1)}{(\sqrt{3} - 1)(\sqrt{3} + 1)} = \frac{2(\sqrt{3} + 1)}{(\sqrt{3})^2 - 1^2} = \frac{2(\sqrt{3} + 1)}{3 - 1} = \sqrt{3} + 1
\]
\subsubsection*{\color{myred} Application}
Soit à rendre rationnels les dénominateurs des rationnels suivants :
\[
I = \frac{2 - \sqrt{3}}{4 + \sqrt{7}} \quad ; \quad J = \frac{1 + \sqrt{5}}{\sqrt{3} - 6} \quad ; \quad K = \frac{7}{\sqrt{6} - \sqrt{5}}
\]
\subsection*{C. Comparaison de réels comportant des radicaux}
\subsubsection*{\color{myred} 1) Activité}
Comparer $2\sqrt{3}$ et $3\sqrt{2}$ en comparant leurs carrés.

\subsubsection*{\color{myred} 2) Propriété et Exemple}
Pour tous nombres réels positifs $a$ et $b$ : 
\[
a < b \iff a^2 < b^2 \iff \sqrt{a} < \sqrt{b}
\]
\textit{Exemple :} 

$(2\sqrt{3})^2 = 4 \times 3 = 12$ et $(3\sqrt{2})^2 = 9 \times 2 = 18$. 

Comme $12 < 18$, alors $2\sqrt{3} < 3\sqrt{2}$.

\subsubsection*{\color{myred} Exemple}

Simplifier et ranger par ordre croissant : 
$A = 3\sqrt{8}$ ; $B = 2\sqrt{18}$ ; $C = 5\sqrt{2}$.
\subsubsection*{\color{myred} Exercice d’application}
Soient : $a = 7 + 2\sqrt{2}$ \quad et \quad $b = 7 - 2\sqrt{2}$.
\begin{enumerate}[label=\textbf{\arabic*)}]
    \item Calcule $\frac{a}{b}$.
    \item Sachant que $1{,}414 < \sqrt{2} < 1{,}415$, encadre $\frac{a}{b}$.
    \item Montre que $a + b$ et $a^2 + b^2$ sont des entiers.
\end{enumerate}
%=========================================================
%    V. Valeur absolue d’un réel
%=========================================================
\section*{\color{myred}V. Valeur absolue d’un réel}

\subsection*{\color{myred} 1) Activité}
Quelle est la distance entre $0$ et $5$ sur une droite graduée ? Entre $0$ et $-5$ ?

\subsection*{\color{myred} 2) Propriété}
La valeur absolue d'un nombre réel $x$, notée $|x|$, représente sa distance à zéro. Elle est toujours positive.
\[
|x| = \begin{cases} x & \text{si } x \ge 0 \\ -x & \text{si } x < 0 \end{cases}
\]
\textit{Exemples :} 

$|-7| = 7$

$|\sqrt{3} - 2| = 2 - \sqrt{3}$ (car $\sqrt{3} < 2$).

%=========================================================
%    VI. Racine carrée du carré d’un réel
%=========================================================
\section*{\color{myred}VI. Racine carrée du carré d’un réel}

\subsection*{\color{myred} 1) Activité}
Calculer $\sqrt{5^2}$ et $\sqrt{(-5)^2}$. Que constate-t-on ?

\subsection*{\color{myred} 2) Propriété}
Pour tout nombre réel $x$, on a :
$\sqrt{x^2} = |x|$

\textbf{Remarques :} 
\begin{itemize}
    \item Si $x \ge 0$, alors $\sqrt{x^2} = x$.
    \item Si $x < 0$, alors $\sqrt{x^2} = -x$ (qui est un nombre positif).
\end{itemize}

\noindent \textit{Exemple type d'examen :}

Simplifier $E = \sqrt{(1 - \sqrt{2})^2}$.

\textbf{Solution :} 

$\sqrt{(1 - \sqrt{2})^2} = |1 - \sqrt{2}|$. Comme $1 < \sqrt{2}$, $1 - \sqrt{2} < 0$, donc $|1 - \sqrt{2}| = -(1 - \sqrt{2}) = \mathbf{\sqrt{2} - 1}$.
\subsection*{\color{myred} 3) Exercice d'application}

Écris sans radical les nombres réels suivants :
\[
\sqrt{(-30)^2} \quad ; \quad \sqrt{(5 - \sqrt{6})^2} \quad ; \quad \sqrt{(2x - 1)^2} \quad ; \quad \sqrt{(3 + \sqrt{5})^2} \quad ; \quad \sqrt{(3 - \sqrt{7})^2}
\]

\vspace{0.3cm}
\hrule
\vspace{0.3cm}

\subsection*{\color{myred} Correction}

Pour tout nombre réel $a$, on utilise la propriété : $\sqrt{a^2} = |a|$.

\begin{itemize}
    \item $\mathbf{\sqrt{(-30)^2}} = |-30| = \mathbf{30}$
    
    \item $\mathbf{\sqrt{(5 - \sqrt{6})^2}} = |5 - \sqrt{6}|$ \\
    Comme $5 = \sqrt{25}$ et $\sqrt{25} > \sqrt{6}$, on a $5 - \sqrt{6} > 0$. \\
    Donc : $|5 - \sqrt{6}| = \mathbf{5 - \sqrt{6}}$
    
    \item $\mathbf{\sqrt{(2x - 1)^2}} = \mathbf{|2x - 1|}$ \\
    \textit{Remarque :} $|2x - 1| = \begin{cases} 2x - 1 & \text{si } x \ge \frac{1}{2} \\ -(2x - 1) = 1 - 2x & \text{si } x < \frac{1}{2} \end{cases}$
    
    \item $\mathbf{\sqrt{(3 + \sqrt{5})^2}} = |3 + \sqrt{5}|$ \\
    Comme $3 + \sqrt{5} > 0$, on a : $|3 + \sqrt{5}| = \mathbf{3 + \sqrt{5}}$
    
    \item $\mathbf{\sqrt{(3 - \sqrt{7})^2}} = |3 - \sqrt{7}|$ \\
    Comme $3 = \sqrt{9}$ et $\sqrt{9} > \sqrt{7}$, on a $3 - \sqrt{7} > 0$. \\
    Donc : $|3 - \sqrt{7}| = \mathbf{3 - \sqrt{7}}$
\end{itemize}
\end{document}